% Part: sets-functions-relations
% Chapter: sets
% Section: russells-paradox

\documentclass[../../../include/open-logic-section]{subfiles}


\begin{document}

\olfileid[pt]{sfr}{set}{rus}
\olsection{Paradoxo de Russell}

A extensionalidade autoriza a notação $\Setabs{x}{\phi(x)}$, para
\emph{o} conjunto dos $x$ tais que~$\phi(x)$. Contudo, tudo o que a
extensionalidade \emph{realmente} autoriza é o seguinte pensamento.
\emph{Se} há um conjunto cujos membros são exatamente os que satisfazem $\phi$,
\emph{então} há apenas um tal conjunto. Dito de outro modo: fixada
alguma~$\phi$, o conjunto $\Setabs{x}{\phi(x)}$ é único, \emph{se existir}.

Mas este condicional é importante!{} Crucialmente, nem toda propriedade
se presta à \emph{compreensão}. Isto é, algumas propriedades \emph{não}
definem conjuntos. Se todas definissem, cairíamos em
contradições flagrantes. O exemplo mais famoso disso é o Paradoxo de Russell.

Os conjuntos podem ser !!{element}s de outros conjuntos---por exemplo, o conjunto potência
de um conjunto~$A$ é constituído por conjuntos. Por isso faz sentido perguntar ou
investigar se um conjunto é !!a{element} de outro conjunto. Pode um conjunto ser
membro de si próprio?  Nada na ideia de conjunto parece excluir
isso. Por exemplo, se \emph{todos} os conjuntos formam uma coleção de
objetos, alguém poderia pensar que podem reunir-se num único
conjunto---o conjunto de todos os conjuntos. E ele, sendo conjunto, seria !!a{element}
do conjunto de todos os conjuntos. 

O Paradoxo de Russell surge quando consideramos a propriedade de não ter
a si próprio como !!a{element}, de ser \emph{não auto-membro}. E se
supusermos que há um conjunto de todos os conjuntos que não têm a si próprios como
!!a{element}? Será que
\[
R = \Setabs{x}{x \notin x}
\]
existe? Verifica-se que podemos provar que não.

\begin{thm}[Paradoxo de Russell]\ollabel{thm:russells-paradox}
	Não há conjunto $R = \Setabs{x}{x \notin x}$.
\end{thm}

\begin{proof}
Se $R = \Setabs{x}{x \notin x}$ existisse, então
$R \in R$ se e só se $R \notin R$, o que é uma contradição.
\end{proof}

\begin{tagblock}{novice}
\begin{explain}
Vamos percorrer esta prova mais devagar. Se $R$ existe, faz sentido perguntar se $R \in
R$ ou não. Suponhamos que de fato $R \in R$. Agora, $R$~foi definido como o conjunto de todos os
conjuntos que não são !!{element}s de si próprios. Logo, se $R \in R$,
então $R$ não tem a propriedade definidora de $R$. Mas só os conjuntos
que têm essa propriedade estão em~$R$, logo $R$ não pode ser !!a{element}
de~$R$, isto é, $R \notin R$. Mas $R$ não pode ao mesmo tempo ser e não ser
!!a{element} de~$R$, logo temos uma contradição.

Como a suposição de que $R \in R$ leva a uma contradição, temos
$R \notin R$. Mas isto também leva a uma contradição!{} Pois se $R
\notin R$, então $R$ tem a propriedade definidora de $R$, e assim $R$ seria
!!a{element} de $R$, tal como todos os outros conjuntos não auto-membros.
E de novo, $R$ não pode ao mesmo tempo não ser e ser !!a{element} de~$R$.
\end{explain}
\end{tagblock}

\begin{digress}
Como montar uma teoria dos conjuntos que evite cair no
Paradoxo de Russell, isto é, que evite afirmar de modo \emph{inconsistente}
que $R = \Setabs{x}{x \notin x}$ existe? Teríamos de estabelecer axiomas que nos
dêem condições muito precisas para dizer quando
os conjuntos existem (e quando não existem). 
	
A teoria dos conjuntos esboçada neste capítulo não faz isso. É
\emph{genuinamente ingênua}. Diz apenas que os conjuntos obedecem à
extensionalidade e que, se temos alguns conjuntos, podemos formar a sua
união, interseção, etc. É possível desenvolver a teoria dos conjuntos com mais
rigor do que isto. \oliflabeldef{cumul:::part}{Esse rigor ficará
reservado à Parte~\olref[cumul][][]{part}. Por enquanto, seguiremos de modo
ingênuo, e tentaremos cuidadosamente evitar contradições.}{}
\end{digress}


\end{document}
